\documentclass[11pt]{article}

\usepackage[a4paper,margin=2.3cm]{geometry}
\usepackage[T1]{fontenc}
\usepackage[utf8]{inputenc}
\usepackage{lmodern}
\usepackage{microtype}
\usepackage{amsmath,amssymb}
\usepackage{booktabs}
\usepackage{array}
\usepackage{graphicx}
\usepackage{siunitx}
\usepackage[authoryear,round]{natbib}
\usepackage{hyperref}
\usepackage{xcolor}
\usepackage{enumitem}
\usepackage{caption}

\hypersetup{
  colorlinks=true,
  linkcolor=black,
  citecolor=black,
  urlcolor=blue
}

\sisetup{detect-all=true}
\setlength{\parindent}{0pt}
\setlength{\parskip}{0.55em}

\newcommand{\figplaceholder}[2]{%
  \IfFileExists{#1}{%
    \includegraphics[width=#2]{#1}%
  }{%
    \fbox{\parbox[c][5.2cm][c]{0.92\linewidth}{\centering
      Figure file not yet committed:\\[0.4em]\texttt{\detokenize{#1}}}}
  }%
}

\title{\textbf{Physics-Guided Soft Sensing under Sparse Laboratory Measurements and Temporal Regime Shifts in Industrial Monoammonium Phosphate Production}}

\author{%
German A. Karnup\textsuperscript{a,b,*}\quad
Oleg A. Telminov\textsuperscript{b}\quad
Andrey N. Sychev\textsuperscript{c}\quad
Evgeny S. Gornev\textsuperscript{b}\\[0.35em]
\small
\href{https://orcid.org/0000-0001-6517-4712}{Karnup ORCID 0000-0001-6517-4712}\quad
\href{https://orcid.org/0000-0002-2358-3689}{Telminov ORCID 0000-0002-2358-3689}
}

\date{}

\begin{document}
\maketitle

\begin{center}
\small
\textsuperscript{a}Moscow Institute of Physics and Technology (National Research University), Dolgoprudny, Moscow Region, Russian Federation\\
\textsuperscript{b}Molecular Electronics Research Institute, JSC (NIIME), Zelenograd, Moscow, Russian Federation\\
\textsuperscript{c}UralAItech LLC, Russian Federation\\
\textsuperscript{*}Corresponding author: \href{mailto:gkarnup@niime.ru}{gkarnup@niime.ru}
\end{center}

\section*{Highlights}
\begin{itemize}[leftmargin=1.5em]
\item Physics-guided soft sensing was evaluated on long-horizon industrial data.
\item Ridge and LSBoost transferred better across temporal shifts.
\item Independent future-period validation exposed weak physics-guided transfer.
\item Deployment requires drift monitoring, recalibration, and fallback logic.
\end{itemize}

\begin{abstract}
Industrial plants collect high-frequency telemetry while critical quality variables may be measured only every few hours. We study soft sensing for an industrial monoammonium phosphate process using 260,642 one-minute records and 1,325 paired laboratory observations. The grey-box model reconstructs latent acid properties from conductivity and combines them with material-balance calculations to estimate molar ratio and downstream density. Validation is entirely chronological: nine pre-gap rolling tests, a short-term regime-transfer test, a strict external post-gap test on 249 future laboratory pairs, and four post-gap rolling comparisons. Process-specific validity filtering is applied before model fitting and temporal alignment. Physics-guided structure improves interpretability but does not ensure temporal robustness. Across regime change and the post-gap period, Ridge and LSBoost remain substantially more accurate than the physics-guided model, demonstrating the need for regime-aware updating of industrial virtual telemetry.
\end{abstract}

\noindent\textbf{Keywords:} soft sensor; virtual sensor; hybrid modeling; process monitoring; regime shift; monoammonium phosphate

\section*{Nomenclature}

\begin{tabular}{@{}lll@{}}
\toprule
Symbol & Definition & Unit\\
\midrule
$c$ & preprocessed conductivity signal used as a proxy for acid state & process unit\\
$Q_A$ & clarified phosphoric-acid volumetric flow & m$^3$ h$^{-1}$\\
$\dot m_N$ & ammonia mass flow & kg h$^{-1}$\\
$Q_W$ & process-water volumetric flow & m$^3$ h$^{-1}$\\
$w_{P_2O_5}$ & mass fraction of $P_2O_5$ in the acid stream & fraction or \%\\
$\rho_A$ & phosphoric-acid density & kg m$^{-3}$\\
$R$ & $NH_3/H_3PO_4$ molar ratio & dimensionless\\
$\rho_M$ & downstream solution density & g cm$^{-3}$\\
$\tau$ & time lag between telemetry and laboratory reference & min\\
$W$ & aggregation-window length & min\\
\bottomrule
\end{tabular}

\section{Introduction}

Chemical-process plants routinely collect high-frequency online measurements while the variables that most directly describe product quality are measured less frequently by laboratory analysis or manual sampling. Soft sensors address this mismatch by inferring hard-to-measure quality variables from continuously available process measurements \citep{kadlec2009,souza2016}.

Industrial deployment is substantially harder than benchmark soft-sensor development. Real plant data contain missing values, sensor faults, short transients, asynchronous sampling, changing operating conditions, and uncertain reference measurements. Data preparation, temporal alignment, and model-update strategy can therefore be as important as model class selection \citep{kadlec2009,offermans2024,dai2025}.

Hybrid and physics-guided models are often proposed as a way to improve interpretability and extrapolation by embedding empirical relationships inside process knowledge rather than learning an unconstrained input-output map \citep{sansana2021,bradley2022,schweidtmann2024,mousa2025}. Industrial soft sensors are also increasingly treated as a digitalization layer that converts existing sensing infrastructure into higher-level virtual measurements for monitoring and decision support \citep{hamid2022,pietrasik2024,boskabadi2025}.

Fertilizer production provides a relevant test case because composition and neutralization variables strongly affect product quality while several feed and intermediate-state properties are not continuously available. Hybrid soft sensing has previously been demonstrated for nutrient-content estimation in compound-fertilizer production \citep{fu2007}. The novelty claimed here is therefore not the generic application of a soft sensor to fertilizer production. Instead, this study examines an industrial monoammonium phosphate (MAP) process with four characteristics that are important for real deployment: one-minute telemetry versus laboratory measurements every few hours, latent feed properties, substantial data-quality problems, and a pronounced operating-regime shift during the observation interval.

This study focuses on the temporal robustness of physics-guided soft sensing under sparse laboratory supervision: \emph{how reliably does the model generalize under chronological hold-out testing, regime change, and an independent post-gap period, and how does its performance compare with simpler data-driven baselines?}

The contributions are:
\begin{enumerate}[leftmargin=1.8em]
\item a real industrial multi-rate soft-sensing case with 260,642 one-minute records and 1,325 paired laboratory observations;
\item a reproducible chronological evaluation using fixed 14-day training and 7-day test windows;
\item comparison of physics-guided, Ridge, physics-only, and LSBoost models on identical temporal splits;
\item explicit separation of process-validity failures from genuine model-transfer limitations;
\item a strict external post-gap validation showing the effect of long-term temporal shift.
\end{enumerate}

\section{Industrial process and data}

\subsection{Process description}

The investigated process is part of an industrial MAP production line. Clarified phosphoric acid is transferred through an intermediate stage to a tubular neutralization reactor. Ammonia is introduced for neutralization and process water is adjusted to influence downstream solution density. The principal reaction is represented in simplified form as
\begin{equation}
NH_3 + H_3PO_4 \rightarrow NH_4H_2PO_4.
\end{equation}

The two quality variables considered in this work are the $NH_3/H_3PO_4$ molar ratio $R$ and downstream solution density $\rho_M$. The nominal operating targets considered in the study were
\begin{equation}
R_{\mathrm{target}}=1.06,
\qquad
\rho_{M,\mathrm{target}}=1.21\ \mathrm{g\,cm^{-3}}.
\end{equation}

The reference values were obtained from sparse laboratory/operator measurements rather than at the one-minute frequency of the process historian. The average interval between laboratory measurements was approximately 4~h. Two phosphoric-acid properties needed by the material-balance calculation, the $P_2O_5$ fraction and acid density, were not continuously measured with sufficient reliability and were therefore treated as latent variables.

\subsection{Data sources}

The full study covers 1 March--29 September 2026 and preserves the one-month data gap from 6 June to 7 July.

\begin{table}[htbp]
\centering
\caption{Data sources used in the study.}
\label{tab:data}
\begin{tabular}{@{}p{3.0cm}p{4.3cm}p{4.4cm}p{2.7cm}@{}}
\toprule
Data source & Variables & Sampling / count & Role\\
\midrule
Plant telemetry & conductivity, acid flow, ammonia flow, water flow & 260,642 one-minute records; 184,387 jointly valid after preprocessing & model inputs\\
Laboratory/operator references & molar ratio and solution density & 1,325 paired observations & calibration and evaluation targets\\
\bottomrule
\end{tabular}
\end{table}

The dataset comprises 139,679 pre-gap process rows and 120,961 post-gap rows. Laboratory coverage comprises 860 paired observations before the gap and 465 paired observations after the gap. For the frozen pre-gap configuration, 760 pre-gap and 249 post-gap laboratory pairs are aligned; for the adapted post-gap configuration, 256 post-gap pairs are aligned.

\subsection{Data quality and preprocessing}

The source data contain noisy process signals, short spikes and dropouts, missing or invalid values, incomplete time-series segments, and operating-state changes. 

The preprocessing workflow is:
\begin{enumerate}[leftmargin=1.8em]
\item read the native one-minute process series;
\item sort channels jointly by timestamp and preserve large acquisition gaps;
\item select the active redundant flow channel where applicable;
\item apply process-specific validity rules to exclude invalid and non-production measurements;
\item apply the authoritative conductivity filter and flow despiking;
\item apply process plausibility and range checks;
\item construct timestamps at which all required inputs are jointly valid;
\item associate each laboratory observation with an eligible process window determined by lag $\tau$ and window length $W$.
\end{enumerate}

Process-specific validity rules were applied to the telemetry channels to remove invalid measurements, short abnormal excursions, and observations outside the production domain before smoothing, model fitting, and laboratory-window aggregation. The same preprocessing rules were applied consistently across the analyzed periods and did not use laboratory target values. Because the one-month data gap exceeds the continuity limit, temporal filters reset at gaps longer than 5 minutes; the gap-aware processing therefore produces two independent telemetry segments.

For retrospective validation, the evaluated preprocessing pipeline retains centered filtering and signal-reconstruction operations, including robust windows, Savitzky--Golay smoothing, and PCHIP repair. These operations can use nearby future process values. Laboratory alignment also uses a centered averaging window around $t_{\mathrm{lab}}-\tau$; when $\tau<W/2$, the window can extend beyond the nominal laboratory timestamp. No outer-test laboratory targets are used in fitting or hyperparameter selection, but a prospective online implementation should replace these process-data operations with causal alternatives.

\subsection{Operating regimes and data gap}

The conductivity trajectory changes abruptly on 21 May 2026. We define Regime~A as the period before 21 May and Regime~B as the period beginning 21 May. The boundary is used as an externally observed process change, not as a label learned by the prediction models. A one-month data gap from 6 June to 7 July is preserved explicitly and separates the earlier and later observation periods; no cause for this interruption is assumed in the analysis. The complete temporal coverage of process telemetry and sparse laboratory references across both periods is summarized in Fig.~\ref{fig:extendedcoverage}.

\begin{figure}[htbp]
\centering
\figplaceholder{figures/fig_extended_validation_timeline.pdf}{0.98\linewidth}
\caption{Extended process and laboratory coverage across the one-month data gap. The upper panel shows the conductivity signal; the lower panels show the sparse laboratory molar-ratio and density references used for calibration and validation.}
\label{fig:extendedcoverage}
\end{figure}

\section{Methodology}

The overall soft-sensor architecture is summarized in Fig.~\ref{fig:soft_sensor_architecture}. The main solid-arrow pathway represents online inference. Minute-scale process telemetry first passes process-validity filtering and temporal alignment; the preprocessed conductivity signal prepared for algorithmic use is then used to reconstruct latent acid properties, which are propagated through the physics-guided material-balance model to obtain the virtual molar-ratio and solution-density estimates. The sparse laboratory measurements are not direct runtime inputs. Instead, they are used only for calibration, chronological model selection, and validation, as indicated by the dashed training pathway. This separation between the online inference path and the calibration path is important for interpreting the model as a deployable virtual-sensing layer rather than as an interpolation of contemporaneous laboratory measurements.

\begin{figure}[htbp]
\centering
\figplaceholder{figures/soft_sensor_architecture.pdf}{0.98\linewidth}
\caption{Architecture of the physics-guided soft sensor and its calibration/validation workflow. The conductivity input denotes the preprocessed signal prepared for algorithmic use. Solid arrows denote the online inference pathway, whereas dashed arrows denote the training, calibration, and validation pathway based on sparse laboratory reference measurements.}
\label{fig:soft_sensor_architecture}
\end{figure}

\subsection{Physics-guided soft sensor}

The soft sensor is a grey-box sequence rather than a direct unconstrained regression. Conductivity is first mapped to a latent $P_2O_5$ fraction,